\subsection{Universality of tree-gauge curl coercivity}
\label{app:lambda_star_universal}

The constant $\lambda_\ast$ used in the Lane~C convexity/scaling-gap module is a \emph{purely combinatorial} constant.
It is determined by the discrete curl operator on a single $2\times2\times2\times2$ block with a fixed spanning-tree gauge choice.
The gauge group enters only through a trivial factorization over Lie-algebra components.

\begin{lemma}[Lie-component factorization and rank-independence]
\label{lem:lambda_star_factorization}
Let $\mathcal Q$ be the canonical $2^4$ coarse block.
Let $D_{\mathrm{red}}^{(1)}\in\mathbb Z^{24\times 17}$ denote the reduced scalar curl matrix on $\mathcal Q$ in the canonical tree gauge
(i.e. acting on real-valued $1$-cochains on edges).
Let $G$ be any compact Lie group with Lie algebra $\mathfrak g$ and $d_G:=\dim\mathfrak g$.
Identify $\mathfrak g\cong\mathbb R^{d_G}$ by choosing an orthonormal basis for the fixed bi-invariant inner product.
Then the reduced curl map on $\mathfrak g$-valued $1$-cochains factorizes as
\begin{equation}
\label{eq:Dred_factorization}
D_{\mathrm{red}}\ =\ D_{\mathrm{red}}^{(1)}\otimes I_{d_G},
\qquad
D_{\mathrm{red}}^{\ast}D_{\mathrm{red}}\ =\ Q\otimes I_{d_G},
\qquad
Q:=(D_{\mathrm{red}}^{(1)})^{\ast}D_{\mathrm{red}}^{(1)}\in\mathbb Z^{17\times 17}.
\end{equation}
In particular,
\begin{equation}
\label{eq:lambda_star_universal_identity}
\lambda_{\min}(D_{\mathrm{red}}^{\ast}D_{\mathrm{red}})\ =\ \lambda_{\min}(Q),
\end{equation}
and the quantity on the right-hand side is independent of the family parameter $n$ for $\Spin(2n{+}1)$, $\Sp(2n)$, and $\Spin(2n)$.
\end{lemma}

\begin{proof}
Fix an orthonormal basis $(e_a)_{a=1}^{d_G}$ of $\mathfrak g$.
A $\mathfrak g$-valued edge field $A=(A_{\mathrm{edge}})$ decomposes uniquely as
$A_{\mathrm{edge}}=\sum_{a=1}^{d_G} A^{(a)}_{\mathrm{edge}}\,e_a$ with scalar components $A^{(a)}$.
The discrete curl on a plaquette is defined by an oriented signed sum of edge variables,
\[ (DA)_{p}:=\sum_{\mathrm{edge}\in\partial p}\sigma_{p,\mathrm{edge}}\,A_{\mathrm{edge}},\qquad \sigma_{p,\mathrm{edge}}\in\{-1,0,1\}. \]
Thus $D$ acts componentwise:
$(DA)^{(a)}=D(A^{(a)})$ for each $a$.
Tree gauge removes a fixed set of edge variables by setting them to zero, which again acts componentwise on the decomposition.
Therefore, in the above basis, the matrix of $D_{\mathrm{red}}$ is the Kronecker product $D_{\mathrm{red}}^{(1)}\otimes I_{d_G}$.
Taking adjoints with respect to the product Euclidean inner product on components yields \eqref{eq:Dred_factorization}.
The identity \eqref{eq:lambda_star_universal_identity} follows from the spectrum of a Kronecker product.
\end{proof}


\subsection{Explicit \texorpdfstring{$Q\to P$}{Q->P} link for the \texorpdfstring{$2^4$}{16} block}
\label{app:lambda_star_Q_to_P}

We now record the explicit reduced curl matrix $Q=(D_{\mathrm{red}}^{(1)})^{\ast}D_{\mathrm{red}}^{(1)}$ for the canonical tree gauge on the $2^4$ block,
compute its characteristic polynomial, and isolate the factor $P$ whose smallest real root equals $\lambda_{\min}(Q)$.

\paragraph{Canonical tree gauge on the $2^4$ block.}
Let the vertex set be $V=\{0,1\}^4$ with the root $0=(0,0,0,0)$.
For each direction $\mu\in\{1,2,3,4\}$ and each vertex $v\in V$ with $v_\mu=0$, let $e=(v,\mu)$ denote the oriented edge from $v$ to $v+e_\mu$.
This gives $|E|=32$ edges.
For each plaquette $(v;\mu,\nu)$ with $1\le\mu<\nu\le4$ and $v_\mu=v_\nu=0$, define the scalar abelianized curl
\[
(DA)_{(v;\mu,\nu)} \,=\, A_{(v,\mu)} + A_{(v+e_\mu,\nu)} - A_{(v+e_\nu,\mu)} - A_{(v,\nu)}.
\]
Tree gauge fixes the $15$ edges of the spanning tree obtained by assigning to each non-root vertex $v\neq 0$ the parent edge
$(v-e_{\mu(v)},\mu(v))$, where $\mu(v):=\max\{\mu: v_\mu=1\}$.
Removing these columns yields $D_{\mathrm{red}}^{(1)}\in\mathbb Z^{24\times 17}$ and
\begin{equation}
\label{eq:Q_def}
Q:=(D_{\mathrm{red}}^{(1)})^{\ast}D_{\mathrm{red}}^{(1)}\in\mathbb Z^{17\times17}.
\end{equation}

\paragraph{Explicit matrix $Q$.}
In the ordered basis of the $17$ unfixed edges induced by the above convention, $Q$ is the following symmetric integer matrix:
{\tiny
\setcounter{MaxMatrixCols}{20}
\[
Q=\begin{pmatrix}
3 & -1 & -1 & 0 & 0 & -1 & 0 & 0 & -1 & 0 & 0 & 0 & 1 & 1 & 0 & 0 & 0 \\
-1 & 3 & -1 & 0 & 0 & 0 & -1 & 0 & 1 & 1 & 0 & 0 & -1 & 0 & 0 & 0 & 0 \\
-1 & -1 & 3 & 0 & 0 & 1 & 1 & 0 & 0 & -1 & 0 & 0 & 0 & -1 & 0 & 0 & 0 \\
0 & 0 & 0 & 3 & -1 & -1 & 0 & 0 & 0 & 0 & -1 & 0 & 0 & 0 & 1 & 0 & 0 \\
0 & 0 & 0 & -1 & 3 & 0 & -1 & 0 & 0 & 0 & 1 & 0 & 0 & 0 & -1 & 0 & 0 \\
-1 & 0 & 1 & -1 & 0 & 3 & -1 & 0 & 0 & 0 & 0 & -1 & 0 & -1 & 0 & 1 & 0 \\
0 & -1 & 1 & 0 & -1 & -1 & 3 & 0 & 0 & -1 & 0 & 1 & 0 & 0 & 0 & -1 & 0 \\
0 & 0 & 0 & 0 & 0 & 0 & 0 & 3 & -1 & 0 & -1 & 0 & 0 & 0 & 0 & 0 & 0 \\
-1 & 1 & 0 & 0 & 0 & 0 & 0 & -1 & 3 & -1 & 0 & -1 & -1 & 0 & 0 & 0 & 1 \\
0 & 1 & -1 & 0 & 0 & 0 & -1 & 0 & -1 & 3 & 0 & 1 & 0 & 0 & 0 & 0 & -1 \\
0 & 0 & 0 & -1 & 1 & 0 & 0 & -1 & 0 & 0 & 3 & -1 & 0 & 0 & -1 & 0 & 0 \\
0 & 0 & 0 & 0 & 0 & -1 & 1 & 0 & -1 & 1 & -1 & 3 & 0 & 0 & 0 & -1 & -1 \\
1 & -1 & 0 & 0 & 0 & 0 & 0 & 0 & -1 & 0 & 0 & 0 & 3 & -1 & 0 & -1 & 1 \\
1 & 0 & -1 & 0 & 0 & -1 & 0 & 0 & 0 & 0 & 0 & 0 & -1 & 3 & 0 & 1 & -1 \\
0 & 0 & 0 & 1 & -1 & 0 & 0 & 0 & 0 & 0 & -1 & 0 & 0 & 0 & 3 & -1 & 0 \\
0 & 0 & 0 & 0 & 0 & 1 & -1 & 0 & 0 & 0 & 0 & -1 & -1 & 1 & -1 & 3 & -1 \\
0 & 0 & 0 & 0 & 0 & 0 & 0 & 0 & 1 & -1 & 0 & -1 & 1 & -1 & 0 & -1 & 3 \\
\end{pmatrix}.
\]
}

\paragraph{Exact characteristic polynomial factorization.}
A direct determinant expansion yields
\begin{align}
\chi_Q(x):=\det(xI-Q)
&=x^{17}-51x^{16}+1176x^{15}-16240x^{14}+149999x^{13}-980451x^{12} \notag\\
&\quad +4681808x^{11}-16630642x^{10}+44353922x^{9}-89060796x^{8} \notag\\
&\quad +134278630x^{7}-150798152x^{6}+124346128x^{5}-73576724x^{4} \notag\\
&\quad +30116880x^{3}-8019792x^{2}+1237248x-82944.\label{eq:chiQ_expanded}
\end{align}
Moreover, $\chi_Q$ admits the exact factorization
\begin{equation}
\label{eq:chiQ_factor}
\chi_Q(x)
=(x-6)(x-1)^3(x^3-8x^2+16x-6)\,P(x),
\end{equation}
where
\begin{align}
P(x)
&=x^{10}-34x^9+489x^8-3884x^7+18664x^6-55904x^5 \notag\\
&\quad +103722x^4-114714x^3+70100x^2-20928x+2304.\label{eq:P_lambda_star}
\end{align}
The identity \eqref{eq:chiQ_factor} follows by polynomial division of \eqref{eq:chiQ_expanded} by $(x-6)(x-1)^3(x^3-8x^2+16x-6)$,
which yields quotient $P$ and zero remainder.

\paragraph{No competing eigenvalue below $1/4$.}
The roots of $(x-6)(x-1)^3$ are $6$ and $1$, hence $>1/4$.
For the cubic factor $f(x)=x^3-8x^2+16x-6$, note that for $x\in[0,\tfrac14]$,
\[
 f'(x)=3x^2-16x+16 \;\ge\; 16-16\cdot\tfrac14 \;=\;12 \;>\;0,
\]
so $f$ is strictly increasing on $[0,\tfrac14]$.
Since $f(\tfrac14)=\tfrac1{64}-\tfrac12+4-6=-\tfrac{159}{64}<0$, it follows that $f(x)<0$ on $[0,\tfrac14]$ and thus $f$ has no real roots there.
Therefore every eigenvalue of $Q$ in $[0,\tfrac14]$ would have to be a root of $P$.


\subsection{Sturm--Habicht root isolation for \texorpdfstring{$P$}{P}}
\label{app:lambda_star_certificate}

We now show that $P$ has no real roots in $[0,\tfrac14]$.
Combined with the previous subsection, this implies $\lambda_{\min}(Q)>\tfrac14$.

\paragraph{Sturm--Habicht (subresultant) sequence.}
Let $P'(x)$ denote the derivative of $P(x)$.
We use the integer Sturm--Habicht sequence (equivalently: the signed subresultant polynomial remainder sequence)
$(S_0,S_1,\dots,S_{10})$ associated to the pair $(P,P')$.
This sequence has the \emph{Sturm property}: for any $a<b$ such that $P(a)P(b)\neq 0$, the number of \emph{distinct} real roots of $P$
in $(a,b]$ equals
\begin{equation}
\label{eq:sturm_variations}
N_{(a,b]}\;=\;V(a)-V(b),
\end{equation}
where $V(t)$ is the number of sign variations in the list $(S_0(t),S_1(t),\dots,S_{10}(t))$ after deleting any zero entries.
See, e.g., \cite[Chapter~2]{BasuPollackRoy2006}.

For completeness, the explicit sequence is:
{\tiny
\begin{align*}
S_{0}(x) &= x^{10} - 34 x^{9} + 489 x^{8} - 3884 x^{7} + 18664 x^{6} - 55904 x^{5} + 103722 x^{4} - 114714 x^{3} + 70100 x^{2} - 20928 x + 2304,\\
S_{1}(x) &= 10 x^{9} - 306 x^{8} + 3912 x^{7} - 27188 x^{6} + 111984 x^{5} - 279520 x^{4} + 414888 x^{3} - 344142 x^{2} + 140200 x - 20928,\\
S_{2}(x) &= - 624 x^{8} + 16488 x^{7} - 177832 x^{6} + 1012256 x^{5} - 3280360 x^{4} + 6076212 x^{3} - 6092828 x^{2} + 2883280 x - 481152,\\
S_{3}(x) &= - 161760 x^{7} + 3651360 x^{6} - 32490048 x^{5} + 145764864 x^{4} - 348422064 x^{3} + 427945008 x^{2} - 235616832 x + 43918848,\\
S_{4}(x) &= 116558720 x^{6} - 2192025280 x^{5} + 15372758720 x^{4} - 50384110080 x^{3} + 78784844800 x^{2} - 52804659200 x + 11503288320,\\
S_{5}(x) &= 156961382528 x^{5} - 2343720929664 x^{4} + 12607138367744 x^{3} - 29677341545728 x^{2} + 29265941425152 x - 8416820994048,\\
S_{6}(x) &= - 1327969697319168 x^{4} + 14372817396860288 x^{3} - 51377384424462976 x^{2} + 68267679340228096 x - 23053088273977344,\\
S_{7}(x) &= - 31880501914859337088 x^{3} + 239230025750838627200 x^{2} - 473271373988031179264 x + 198883701046191943680,\\
S_{8}(x) &= 814010205748396620519168 x^{2} - 3141624232683282745586688 x + 2561510129798987060207616,\\
S_{9}(x) &= 49205066212446239571161554944 x - 108703453356939315792064610304,\\
S_{10}(x) &= -1484006905322599188659461240651776.
\end{align*}
}

\paragraph{Sign table and variation counts.}
For each $k$ set $d_k:=\deg S_k$ and define the integers
\begin{equation}
\label{eq:sturm_integer_defs}
A_{k,0}:=S_k(0),\qquad 
A_{k,1}:=4^{d_k}S_k\!\bigl(\tfrac14\bigr),\qquad 
A_{k,2}:=50^{d_k}S_k\!\bigl(\tfrac{13}{50}\bigr).
\end{equation}
Since each $S_k\in\mathbb Z[x]$, the values $A_{k,0},A_{k,1},A_{k,2}$ are integers; the powers of $4$ and $50$ clear denominators. Each $A_{k,j}$ is obtained by substituting $x=0,\,\tfrac14,\,\tfrac{13}{50}$ into the explicitly displayed polynomial $S_k(x)$ above and multiplying by $4^{d_k}$ or $50^{d_k}$ to clear denominators.
Moreover $4^{d_k},50^{d_k}>0$, hence $\mathrm{sign}\,S_k(\tfrac14)=\mathrm{sign}\,A_{k,1}$ and
$\mathrm{sign}\,S_k(\tfrac{13}{50})=\mathrm{sign}\,A_{k,2}$.
Table~\ref{tab:sturm_habicht_integer_evals} records these integers (with spaces inserted every three digits for readability).
In particular, none of the listed values is zero, so no deletion of zero entries is needed in the variation count.

\begin{table}[t]
\centering
\renewcommand{\arraystretch}{1.05}
\setlength{\tabcolsep}{3pt}
\scriptsize
\begin{tabular}{r r >{\raggedright\ttfamily}p{0.12\linewidth} >{\raggedright\ttfamily}p{0.27\linewidth} >{\raggedright\ttfamily\arraybackslash}p{0.36\linewidth}}
\toprule
$k$ & $d_k$ & $A_{k,0}$ & $A_{k,1}=4^{d_k}S_k(\tfrac14)$ & $A_{k,2}=50^{d_k}S_k(\tfrac{13}{50})$ \\
\midrule
0 & 10 & 2 304 & 16 505 609 & -166 658 660 387 519 751 \\
1 & 9 & -20 928 & -496 282 942 & -3 125 596 796 325 277 570 \\
2 & 8 & -481 152 & -3 805 068 112 & -1 969 541 134 507 775 104 \\
3 & 7 & 43 918 848 & 112 323 055 776 & 4 759 617 805 973 242 080 \\
4 & 6 & 11 503 288 320 & 10 227 147 758 720 & 35 657 229 918 349 172 480 \\
5 & 5 & -8 416 820 994 048 & -2 833 597 260 445 056 & -813 376 624 526 073 786 496 \\
6 & 4 & -23 053 088 273 977 344 & -2 298 333 971 264 887 552 & -53 312 841 855 303 381 480 448 \\
7 & 3 & 198 883 701 046 191 943 680 & 6 081 254 484 236 280 699 008 & 11 430 595 231 050 620 070 177 664 \\
8 & 2 & 2 561 510 129 798 987 060 207 616 & 29 231 675 351 799 058 601 494 272 & 4 499 287 298 024 812 894 755 432 192 \\
9 & 1 & -108 703 453 356 939 315 792 064 610 304 & -385 608 747 215 311 023 597 096 886 272 & -4 795 506 807 085 164 675 178 130 300 928 \\
10 & 0 & -1 484 006 905 322 599 188 659 461 240 651 776 & -1 484 006 905 322 599 188 659 461 240 651 776 & -1 484 006 905 322 599 188 659 461 240 651 776 \\
\bottomrule
\end{tabular}
\caption{Exact integer evaluations used for the Sturm--Habicht variation count. Spaces in the displayed integers are for readability.}
\label{tab:sturm_habicht_integer_evals}
\end{table}

The resulting sign pattern at the three test points is:

\begin{center}
\begin{tabular}{rcccc}
\toprule
$k$ & $\deg S_k$ & $\mathrm{sign}\,S_k(0)$ & $\mathrm{sign}\,S_k(\tfrac14)$ & $\mathrm{sign}\,S_k(\tfrac{13}{50})$ \\
\midrule
0  & 10 & $+$ & $+$ & $-$ \\
1  & 9  & $-$ & $-$ & $-$ \\
2  & 8  & $-$ & $-$ & $-$ \\
3  & 7  & $+$ & $+$ & $+$ \\
4  & 6  & $+$ & $+$ & $+$ \\
5  & 5  & $-$ & $-$ & $-$ \\
6  & 4  & $-$ & $-$ & $-$ \\
7  & 3  & $+$ & $+$ & $+$ \\
8  & 2  & $+$ & $+$ & $+$ \\
9  & 1  & $-$ & $-$ & $-$ \\
10 & 0  & $-$ & $-$ & $-$ \\
\bottomrule
\end{tabular}
\end{center}
Thus
\[
V(0)=5,\qquad V\bigl(\tfrac14\bigr)=5,\qquad V\bigl(\tfrac{13}{50}\bigr)=4.
\]
By \eqref{eq:sturm_variations}, $P$ has no real roots in $(0,\tfrac14]$ and exactly one real root in $(\tfrac14,\tfrac{13}{50}]$.
Since $P(0)>0$, this implies $P$ has no real roots in $[0,\tfrac14]$.

\paragraph{Conclusion.}\label{concl:lambda_star_gt_quarter}
The smallest real root of $P$ satisfies $\lambda_\ast>\tfrac14$.
Consequently, by \eqref{eq:lambda_star_universal_identity}, the same bound holds for the full operator $D_{\mathrm{red}}^{\ast}D_{\mathrm{red}}$
for every compact gauge group $G$ and in particular for $\Spin(2n{+}1)$ ($n\ge2$), $\Sp(2n)$ ($n\ge3$), and $\Spin(2n)$ ($n\ge4$).
